\begin{table*}[!t]
\centering
\footnotesize
\renewcommand{\arraystretch}{1}  % add vertical breathing room
\setlength{\tabcolsep}{9pt}
\captionsetup{font=small}

\caption{\textbf{Annotation Rubric for Speech Naturalness.}
Each dialog is evaluated along multiple sub-dimensions that capture different aspects of conversational speech quality.}
\vspace{-0.6em}

\begin{tabular}{p{3.2cm} p{12.2cm}}
\toprule
\textbf{Sub-metric} & \textbf{Description} \\
\midrule
\textbf{Expressive Intensity} &
The degree of expressiveness, reflecting how emotionally vivid or subdued the speech sounds. \\

\textbf{Expressive Correctness} &
The appropriateness of the expressiveness w.r.t. the conversational context and semantic content. \\

\textbf{Intonation} &
The quality of intonation, including pitch variation, stress patterns, and overall prosodic naturalness. \\

\textbf{NSVS}&
The naturalness of non-speech vocalizations and filler words (e.g., breaths, pauses, hesitations). \\

\textbf{Mispronunciation} &
Whether mispronunciations are present in the speech, reflecting pronunciation accuracy. \\

\textbf{Pacing} &
The naturalness of the pacing or speaking rate, including timing, pauses, and temporal flow. \\

\midrule
\textbf{Human-likeness} &
An overall rating of how closely the speech resembles natural human conversational speech. \\
\bottomrule
\end{tabular}
\label{table:rubric}
\vspace{-1em}
\end{table*}
